\section{Implementation}
\label{sec:implementation}

\subsection{Simulator and reproducibility}
\label{sec:sim-implementation}

The evaluation harness is a discrete-time simulator written in Python\,3.14 with
NumPy\,2.5 and SciPy\,1.18. It generates a 24-node deployment with the tier mix
of \cref{sec:system-model}, a latent threat state, four context modalities and
per-request access decisions. Each run has 2{,}000 epochs; ten seeds are used for
every scenario, and all ten policies are replayed against the byte-identical
observation and request stream generated for that seed and scenario. Comparisons
are therefore paired, and a policy cannot gain by seeing a different sequence of
requests. The full matrix is 800 policy runs (10 seeds $\times$ 8 scenarios
$\times$ 10 policies) and completes in about 50 seconds on eight cores.

The simulator is \emph{parameterised}, not a replay. The observation model, the
transition matrix and the attack processes are explicit; no raw packet capture is
loaded at run time. The public IoT-23 capture is used offline to calibrate one
prior (the benign-to-malicious network-intensity contrast) and is documented in
\cref{sec:trace}. This distinction is deliberate: a simulator that claims to
replay a dataset but generates random numbers is not reproducible, and a
trace-driven evaluation would answer a different question (can the monitor
recognise known attack signatures?) than the one asked here (does the controller
remain safe when the context is adversarially chosen?).

\subsection{Scenarios and policies}
\label{sec:scenarios}

The eight scenarios share the same node topology and request process; they
differ in the latent state process and the poisoning channel.

\begin{description}
  \item[None.] Benign operation with occasional reconnaissance. This is the
  utility reference, not a safety test.
  \item[DoS.] Three non-overlapping attack windows in which the latent state is
  DoS; the network and security modalities show congestion and inspection
  failures. No context spoofing is applied, so the attack is observable.
  \item[Downgrade.] The same windows, but the adversary moves the network and
  security modalities toward the benign distribution while leaving the attested
  device modality and the cross-node dispersion intact. Physical effects remain,
  but the network evidence is consistent with a calm plant. This is the in-model
  poisoning attack: it assumes the anchor is trustworthy.
  \item[TEE.] An explicit out-of-model stress test in which the adversary also
  spoofs the attested device modality, i.e.\ the node's trusted execution
  environment is assumed subverted. It is the only scenario that instantiates
  the fifth latent state (compromised); the other attack scenarios use DoS,
  Byzantine or no serious state. This scenario exists to separate two claims:
  the invariant guarantee, which still holds, from belief quality, which does
  not. It is not used to argue that the architecture detects a TEE compromise.
  \item[Energy.] The adversary inflates the apparent threat: the network and
  security modalities look like an attack while the plant is benign. The goal is
  denial-of-sleep and over-provisioning.
  \item[Byzantine.] Ten of 24 nodes collude and report fabricated benign context.
  The honest majority still determines the median for most features, but the
  cross-node dispersion rises. This tests aggregation robustness, not majority
  capture.
  \item[Resonant.] The adversary alternates downgrade and energy poisoning at a
  period close to the controller's dwell time, attempting to force posture
  thrashing.
  \item[Adaptive.] The adversary starts with partial spoofing, increases the
  spoofed fraction over the first 600 epochs, and then holds the anomaly score
  just below the public point-estimate threshold (0.55). It uses the published
  policy parameters but does not solve the defender's POMDP.
\end{description}

The ten policies are: Static-High (always the strongest posture, with Ascon at
the AEAD level), Static-Med (a uniform middle posture, raised to each
operation's floor), Static-Low (always the weakest), Threshold-CAAC (a
univariate point-estimate anomaly threshold on the same feature the attacker
manipulates, raised to the structural floor; a multivariate or learned detector
is not evaluated, which is a limitation of the baseline set), Risk-Only
(Bayesian belief, no invariants), Invariant-Only
(invariants, no belief; uses the anomaly point estimate), Full$-$Hyst (full
controller without hysteresis or margin), Full$-$Cons (full controller with
uniform modality fusion), Full (the proposed controller) and Oracle (knows the
latent state). Static-High and Static-Med use the same Ascon primitive as the
adaptive policies, so the comparison is not confounded by cipher choice. The
deliberately unconstrained baselines (Static-Low, Risk-Only) do not install the
node-side PEP re-check; every policy that claims the invariant guarantee does.
Making the unconstrained baselines structurally safe would erase the
counterfactual that RQ1 measures.

\subsection{Analytical cost model}
\label{sec:cost-model}

Latency and energy are estimated from per-layer operation costs rather than
measured on silicon. For action
$a=(a_{\mathrm{auth}},a_{\mathrm{cry}},a_{\mathrm{ids}},a_{\mathrm{rate}},a_{\mathrm{ttl}})$,
\begin{equation}
  \ell(a)=\ell_0+\sum_{\ell}\ell_{\ell}(a_\ell),\qquad
  e(a)=e_0+\sum_{\ell}e_{\ell}(a_\ell),
  \label{eq:cost}
\end{equation}
where the per-layer constants are derived from datasheet-level operation counts
for the three tiers and $\ell_0,e_0$ cover request validation. The tier
multipliers are $1.00$, $0.60$ and $0.35$ for Class-1, Class-2 and gateway
respectively, and the reported latency uses the tier distribution of the
deployment. The resulting numbers are \emph{analytical estimates}; they are
useful for comparing postures and for checking that the controller fits an MCU
budget, but they are not a substitute for a hardware measurement. SRAM and flash
footprints are likewise model estimates: 9--14\,KB SRAM and 58--86\,KB flash
depending on whether the policy needs the belief filter.

The cost model is intentionally coarse. Its role is to make the utility
trade-off non-trivial (stronger actions must buy enough detection to justify
their latency and energy), not to predict a specific board's power draw. A
deployment would replace \eqref{eq:cost} with a measured profile; the controller
only consumes the resulting scalar costs.

\subsection{Verification}
\label{sec:verification}

The invariant layer is verified exhaustively rather than sampled. A unit test
enumerates all 243 actions for each of the four operation classes and checks
that the structural mask matches the independent restatement of I1/I2, that the
conditional mask is a subset of the structural mask for a 41-point grid of
$\attack\in[0,1]$, and that the admissible set is never empty. A second test
runs all adaptive policies on a real stream and asserts that no emitted action
violates I1/I2. Belief normalisation, trust bounds, AUC/TPR helpers, paired
determinism and the trace-calibration consistency are covered by 23 tests, all
passing. These tests do not prove the cryptographic mechanisms; they prove that
the decision layer never selects an action outside the declared envelope.
